% =====================================================================
%  DRAFT FOR REVIEW — 4 missing glossary entries (annex.tex)
%  These four frameworks are columns in the Phase 3 matrix but had no
%  glossary entry yet. Format matches your existing entries
%  (metadata longtable + Objective / In the ecosystem / Status flag).
%  PLACEMENT:
%   - CMO  -> sector-regulation section, next to FCR / IUU.
%   - F-gas, Whistleblowing, Pay Transparency -> EU horizontal / general
%     regulation section, next to OSH / Working Time / Forced Labour / PPWR.
%  Dates verified June 2026; re-check before final citation where flagged.
% =====================================================================

% --------------------------------------------------------------------- CMO
CMO --- Common Market Organisation marketing standards

{\def\LTcaptype{none}
\begin{longtable}[]{@{}
  >{\raggedright\arraybackslash}p{(\linewidth - 2\tabcolsep) * \real{0.2382}}
  >{\raggedright\arraybackslash}p{(\linewidth - 2\tabcolsep) * \real{0.7618}}@{}}
\toprule\noalign{}
\endhead
\bottomrule\noalign{}
\endlastfoot
\textbf{Behind it} & EU (Parliament \& Council) \\
\textbf{Reference / founded} & Regulation (EU) No 1379/2013 (CMO for fishery and aquaculture products) \\
\textbf{Status · role} & EU regulation; mandatory, directly applicable · Asker \\
\textbf{Stakeholder source} & Regulator (sector) \\
\textbf{Official source} & \href{https://eur-lex.europa.eu/eli/reg/2013/1379/oj}{eur-lex.europa.eu --- Reg (EU) 1379/2013} \\
\end{longtable}
}

\textbf{Objective.} Sets common marketing standards and \emph{mandatory consumer information} for fishery and aquaculture products: commercial and scientific species name, production method (caught / caught in freshwater / farmed), catch or production area (FAO zone or Member State), fishing-gear category, and whether the product has been defrosted.

\textbf{In the ecosystem.} Binds the processor directly at labelling and point of sale. It is the sector's mandatory transparency floor on origin and method --- the regulatory backbone behind several traceability and sustainable-sourcing disclosures, even though it is an information duty rather than a sustainability-performance standard.

% --------------------------------------------------------------------- F-gas
F-gas Regulation --- Fluorinated greenhouse gases

{\def\LTcaptype{none}
\begin{longtable}[]{@{}
  >{\raggedright\arraybackslash}p{(\linewidth - 2\tabcolsep) * \real{0.2382}}
  >{\raggedright\arraybackslash}p{(\linewidth - 2\tabcolsep) * \real{0.7618}}@{}}
\toprule\noalign{}
\endhead
\bottomrule\noalign{}
\endlastfoot
\textbf{Behind it} & EU (Parliament \& Council) \\
\textbf{Reference / founded} & Regulation (EU) 2024/573 --- in force March 2024 (repealing Reg 517/2014) \\
\textbf{Status · role} & EU regulation; mandatory, directly applicable · Asker \\
\textbf{Stakeholder source} & Regulator (horizontal) \\
\textbf{Official source} & \href{https://eur-lex.europa.eu/eli/reg/2024/573/oj}{eur-lex.europa.eu --- Reg (EU) 2024/573} \\
\end{longtable}
}

\textbf{Objective.} Controls emissions of fluorinated greenhouse gases through an HFC phase-down (quota system), bans on high-GWP equipment, and direct operator duties: leak checks on equipment above 5 tonnes CO\textsubscript{2}-equivalent, leak-detection systems, recovery and certified handling, five-year record-keeping, and reporting via the F-gas portal (Art.~26).

\textbf{In the ecosystem.} Scope-conditional but acutely sector-relevant: refrigeration and freezing are core to seafood processing and the cold chain. It binds the processor where \emph{fluorinated} (HFC/HFO) refrigerants are used; firms on natural refrigerants (ammonia, CO\textsubscript{2}) fall outside its scope. Its refrigerant-emission data feed Scope~1 GHG (E1).

\textbf{Status flag.} \emph{Applies only to fluorinated-gas users --- governed in Phase~3 by a firm-profile include-switch.}

% --------------------------------------------------------------- Whistleblowing
Whistleblowing Directive --- Protection of persons reporting breaches of EU law

{\def\LTcaptype{none}
\begin{longtable}[]{@{}
  >{\raggedright\arraybackslash}p{(\linewidth - 2\tabcolsep) * \real{0.2382}}
  >{\raggedright\arraybackslash}p{(\linewidth - 2\tabcolsep) * \real{0.7618}}@{}}
\toprule\noalign{}
\endhead
\bottomrule\noalign{}
\endlastfoot
\textbf{Behind it} & EU (Parliament \& Council) \\
\textbf{Reference / founded} & Directive (EU) 2019/1937 --- in force for 50--249-employee entities since 17 Dec 2023 \\
\textbf{Status · role} & EU directive; transposed nationally; mandatory · Asker \\
\textbf{Stakeholder source} & Regulator (horizontal) \\
\textbf{Official source} & \href{https://eur-lex.europa.eu/eli/dir/2019/1937/oj}{eur-lex.europa.eu --- Dir (EU) 2019/1937} \\
\end{longtable}
}

\textbf{Objective.} Requires private legal entities with 50 or more employees to establish confidential \emph{internal reporting channels} and procedures, protect reporters from retaliation, give feedback within set timeframes, and keep records of reports.

\textbf{In the ecosystem.} Binds the processor directly as an employer. It is admitted on the same basis as the OSH Framework Directive --- a binding organisation-level governance \emph{mechanism} on a G/S topic, not a public-disclosure mandate. Its content maps to business-conduct (G1) and the worker grievance / raise-a-concern channels in S1-3 and S2-3.

% ----------------------------------------------------------- Pay Transparency
Pay Transparency Directive --- Equal pay and pay transparency

{\def\LTcaptype{none}
\begin{longtable}[]{@{}
  >{\raggedright\arraybackslash}p{(\linewidth - 2\tabcolsep) * \real{0.2382}}
  >{\raggedright\arraybackslash}p{(\linewidth - 2\tabcolsep) * \real{0.7618}}@{}}
\toprule\noalign{}
\endhead
\bottomrule\noalign{}
\endlastfoot
\textbf{Behind it} & EU (Parliament \& Council) \\
\textbf{Reference / founded} & Directive (EU) 2023/970 --- transposition deadline 7 June 2026 \\
\textbf{Status · role} & EU directive; phased by headcount; mandatory · Asker \\
\textbf{Stakeholder source} & Regulator (horizontal) \\
\textbf{Official source} & \href{https://eur-lex.europa.eu/eli/dir/2023/970/oj}{eur-lex.europa.eu --- Dir (EU) 2023/970} \\
\end{longtable}
}

\textbf{Objective.} Strengthens equal-pay enforcement and introduces mandatory \emph{gender pay-gap reporting}: employers with 250+ employees report annually; 150--249 every three years (first report 2027 on 2026 data); 100--149 from 2031. Triggers a joint pay assessment where an unjustified gap of 5\% or more is found.

\textbf{In the ecosystem.} Binds the processor directly above the headcount trigger. Its mandated datapoint (the gender pay gap) corresponds to ESRS~S1-16, which is \emph{outside} the cohort-defined topic universe --- so its demand is recorded in the out-of-frame register rather than scored on an in-universe row.

\textbf{Status flag.} \emph{Applies at or above 150 employees --- governed in Phase~3 by a firm-profile include-switch.}
